% Reconstructed from the 14 August 2026 FINAL PDF.
\documentclass[11pt]{article}
\usepackage[margin=1.1in]{geometry}
\usepackage{amsmath,amssymb,amsthm}
\usepackage{enumitem}
\usepackage[colorlinks=true,linkcolor=blue,citecolor=blue,urlcolor=blue]{hyperref}

\newcommand{\pP}{\textbf{[P]}}
\newcommand{\pC}{\textbf{[C]}}
\newcommand{\pO}{\textbf{[O]}}

\title{\textbf{The Hehl--Datta Coefficient from Relative-Entropy Positivity\\
(Final Updated Version):\\
Sign, Rigidity, and Pure-Information Magnitude}}
\author{Robert W.\ McGwier\\
\small CTO and Staff Scientist, Cohere Technology Group, Herndon, VA}
\date{August 14, 2026 --- Final Updated Version}

\begin{document}
\maketitle

\begin{abstract}
\noindent
This is the final updated version of Paper V of the Emergent Gravity Program.
It incorporates all subsequent results of the series so that every claim
carries its correct final epistemic label.

Sign \pP{}: relative-entropy positivity forces the contact term negative.
Magnitude by rigidity \pP{}: given the Ryu--Takayanagi normalization and
the modular selection rule, the coefficient is uniquely $-3\kappa/16$.
Magnitude from pure information \pP{}: the Fisher--canonical matching is
closed by the null-plane limit together with the absorbability of the
finite-ball residue.

The multi-digit numerical coefficient of the kinematic residue $I_B$,
the de Sitter domain, the non-linear regime, and UV completion remain
\pO{}. No outdated open labels appear on claims that have been closed.

Why this update is required. The original Paper V left the
pure-information magnitude open because it depended on the then-unevaluated
integral $I_B$. That integral has been shown to be non-vanishing and
absorbable, and the matching has been closed. This version brings Paper V
into full consistency with the rest of the series.
\end{abstract}

\section*{Epistemic Conventions}
\pP{} proved/derived; \pC{} conjectured with support; \pO{} open.
Numerical agreement is not proof. Consensus is not proof.

\section{Final Status of the Claims of This Paper}

\begin{enumerate}[leftmargin=1.6em]
\item Sign of the Hehl--Datta coefficient (relative-entropy positivity). \pP{}
\item Magnitude by rigidity (RT normalization + modular selection rule). \pP{}
\item Magnitude from pure information (null-plane limit + absorbability). \pP{}
\item Resulting coefficient $\mathcal{L}_{\mathrm{HD}}=-3\kappa/16\, J_5\cdot J_5$. \pP{}
\item Multi-digit numerical coefficient of the residue $I_B$. \pO{}
\item de Sitter / cosmological domain. \pO{}
\item Non-linear regime and UV completion. \pO{}
\end{enumerate}

\section{Cross-References to the Rest of the Series}
This paper is consistent with and cites:
\begin{itemize}[leftmargin=1.4em]
\item Paper IV (Final Updated) --- unconditional metric sector, non-vanishing and absorbable residue
\item Paper VI --- evaluation establishing non-vanishing of $I_B$
\item Paper VII --- explicit closure of the Fisher--canonical matching
\item Paper VIII (Version 4) --- order-of-magnitude coefficient ($O(1)$) and remaining analytic requirement
\item Paper IX --- domain of validity (AdS closed, de Sitter open)
\item Final Status document (14 August 2026) --- authoritative label table for the whole program
\end{itemize}

\section{Summary}
The Hehl--Datta coefficient is now determined in sign and in magnitude by
quantum-information principles inside the AdS regime. The only open items
that remain are those that lie outside the reach of the present methods.
All labels in this version reflect that final accounting.

\begin{thebibliography}{99}
\bibitem{I} R.W.\ McGwier, ``Entanglement as the Origin of Spacetime Curvature,'' Zenodo (2026).
\bibitem{IV} R.W.\ McGwier, ``Condition NL for Finite Balls'' (Final Updated), Zenodo (2026).
\bibitem{V} R.W.\ McGwier, ``The Hehl--Datta Coefficient\ldots'' (this work, final update), Zenodo (2026).
\bibitem{VI} R.W.\ McGwier, ``Evaluation of the Universal Kinematic Integral,'' Zenodo (2026).
\bibitem{VII} R.W.\ McGwier, ``Pure-Information Determination of the Hehl--Datta Coefficient,'' Zenodo (2026).
\bibitem{VIII} R.W.\ McGwier, ``Status of the High-Precision Coefficient of $I_B$'' (Version 4), Zenodo (2026).
\bibitem{IX} R.W.\ McGwier, ``Domain of Validity of the Emergent Gravity Program,'' Zenodo (2026).
\bibitem{FS} R.W.\ McGwier, ``Emergent Gravity Program: Final Status of Claims,'' Zenodo (2026).
\bibitem{FGHMV} T.\ Faulkner et al., JHEP 03 (2014) 051.
\bibitem{LV} N.\ Lashkari, M.\ Van Raamsdonk, JHEP 04 (2016) 153.
\bibitem{CTT} H.\ Casini, E.\ Teste, G.\ Torroba, J.\ Phys.\ A 50 (2017) 364001.
\bibitem{Hehl} F.W.\ Hehl et al., Rev.\ Mod.\ Phys.\ 48 (1976) 393.
\end{thebibliography}
\end{document}
